% GENERATED FILE. Edit canonical lessons, then run npm run generate:books.
% canonical-source-hash: fnv1a64:8b35ada9b0ef70da
% canonical-lessons: LA-C251-catulus, LA-C251-tigris, LA-C251-clava, LA-C251-catena, LA-C251-scutum

\chapter{Puppy, Tiger, Club, Chain, Shield}
\label{ch:la-puppy-tiger-club-chain-shield}
\hlchaptermodality{251}

\begin{chapteropening}
\textbf{By the end of this chapter:} \emph{I can say a puppy, a tiger, a club, a chain and a shield.}

Complete the last lesson of chapter 251: I can say a puppy, a tiger, a club, a chain and a shield.
\end{chapteropening}

\section[catulus, catulī]{catulus, catulī — a puppy}

\label{lesson:LA-C251-catulus}



\begin{quote}
Before the new one: say the Latin for a laurel, then the Latin for a violet.
\end{quote}

\subsection*{You'll want to know: catulus, catulī}
\textbf{catulus, catulī} — \textquotedblleft{}a puppy\textquotedblright{}.

\begin{cousinweb}
\textbf{catulus}, \textquotedblleft{}a young animal, a puppy\textquotedblright{}, has no sure origin. Lewis \& Short make it a diminutive of \emph{catus}, \textquotedblleft{}a cat\textquotedblright{}, but \emph{catulus} is far older than that word for cat, so the link is unlikely.
\end{cousinweb}

\begin{grammarlens}[title={What you've built}]
One more word for this chapter.
\end{grammarlens}

\subsection*{Your turn}
\begin{itemize}
\raggedright
  \item \emph{Say it:} \emph{catulus, catulī}
  \item \emph{Say it:} \emph{catulus, catulī}, once more
  \item \emph{Say it:} say \emph{catulus, catulī}
  \item \emph{Recall:} say the Latin for a laurel, then the Latin for a violet, then say \emph{catulus, catulī} again
\end{itemize}

\begin{tcolorbox}[breakable,colback=teal!4,colframe=teal!35!black,title={Before you move on}]
What does catulus, catulī mean? (\textquotedblleft{}A puppy\textquotedblright{}.) Say it once more.
\end{tcolorbox}
\section[tigris, tigris]{tigris, tigris — a tiger}

\label{lesson:LA-C251-tigris}



\begin{quote}
Before the new one: say the Latin for a violet, then the Latin for a puppy.
\end{quote}

\subsection*{You'll want to know: tigris, tigris}
\textbf{tigris, tigris} — \textquotedblleft{}a tiger\textquotedblright{}.

\begin{cousinweb}
\textbf{tigris} was borrowed from Greek \emph{tigris}, \textquotedblleft{}tiger\textquotedblright{} (EtymDB 2.1), which probably came from an Iranian language.
\end{cousinweb}

\begin{grammarlens}[title={What you've built}]
One more word for this chapter.
\end{grammarlens}

\subsection*{Your turn}
\begin{itemize}
\raggedright
  \item \emph{Say it:} \emph{tigris, tigris}
  \item \emph{Say it:} \emph{tigris, tigris}, once more
  \item \emph{Say it:} say \emph{tigris, tigris}
  \item \emph{Recall:} say the Latin for a violet, then the Latin for a puppy, then say \emph{tigris, tigris} again
\end{itemize}

\begin{tcolorbox}[breakable,colback=teal!4,colframe=teal!35!black,title={Before you move on}]
What does tigris, tigris mean? (\textquotedblleft{}A tiger\textquotedblright{}.) Say it once more.
\end{tcolorbox}
\section[clāva, clāvae]{clāva, clāvae — a club}

\label{lesson:LA-C251-clava}



\begin{quote}
Before the new one: say the Latin for a puppy, then the Latin for a tiger.
\end{quote}

\subsection*{You'll want to know: clāva, clāvae}
\textbf{clāva, clāvae} — \textquotedblleft{}a club\textquotedblright{}.

\begin{cousinweb}
\textbf{clāva} is often traced to Proto-Indo-European \emph{*kelh2-}, \textquotedblleft{}to beat, to break\textquotedblright{} (EtymDB 2.1), with \emph{percellō}, \textquotedblleft{}to strike down\textquotedblright{}, in the same family (Lewis \& Short); the link is not certain.
\end{cousinweb}

\begin{grammarlens}[title={What you've built}]
One more word for this chapter.
\end{grammarlens}

\subsection*{Your turn}
\begin{itemize}
\raggedright
  \item \emph{Say it:} \emph{clāva, clāvae}
  \item \emph{Say it:} \emph{clāva, clāvae}, once more
  \item \emph{Say it:} say \emph{clāva, clāvae}
  \item \emph{Recall:} say the Latin for a puppy, then the Latin for a tiger, then say \emph{clāva, clāvae} again
\end{itemize}

\begin{tcolorbox}[breakable,colback=teal!4,colframe=teal!35!black,title={Before you move on}]
What does clāva, clāvae mean? (\textquotedblleft{}A club\textquotedblright{}.) Say it once more.
\end{tcolorbox}
\section[catēna, catēnae]{catēna, catēnae — a chain}

\label{lesson:LA-C251-catena}



\begin{quote}
Before the new one: say the Latin for a tiger, then the Latin for a club.
\end{quote}

\subsection*{You'll want to know: catēna, catēnae}
\textbf{catēna, catēnae} — \textquotedblleft{}a chain\textquotedblright{}.

\begin{cousinweb}
\textbf{catēna} has no sure origin; EtymDB 2.1 records a proposed Proto-Indo-European \emph{*kat-}, \textquotedblleft{}to twist, to plait\textquotedblright{}. English \emph{chain} comes from it through French.
\end{cousinweb}

\begin{grammarlens}[title={What you've built}]
One more word for this chapter.
\end{grammarlens}

\subsection*{Your turn}
\begin{itemize}
\raggedright
  \item \emph{Say it:} \emph{catēna, catēnae}
  \item \emph{Say it:} \emph{catēna, catēnae}, once more
  \item \emph{Say it:} say \emph{catēna, catēnae}
  \item \emph{Recall:} say the Latin for a tiger, then the Latin for a club, then say \emph{catēna, catēnae} again
\end{itemize}

\begin{tcolorbox}[breakable,colback=teal!4,colframe=teal!35!black,title={Before you move on}]
What does catēna, catēnae mean? (\textquotedblleft{}A chain\textquotedblright{}.) Say it once more.
\end{tcolorbox}
\section[scūtum, scūtī]{scūtum, scūtī — a shield}

\label{lesson:LA-C251-scutum}



\begin{quote}
Before the new one: say the Latin for a club, then the Latin for a chain.
\end{quote}

\subsection*{You'll want to know: scūtum, scūtī}
\textbf{scūtum, scūtī} — \textquotedblleft{}a shield\textquotedblright{}.

\begin{cousinweb}
\textbf{scūtum} is probably inherited, from a Proto-Indo-European root \emph{*skei-}, \textquotedblleft{}to cut, to split\textquotedblright{} (EtymDB 2.1): a shield of split boards. Old Irish \emph{sciath}, \textquotedblleft{}shield\textquotedblright{}, is a cousin; Lewis \& Short's link with Greek \emph{skytos}, \textquotedblleft{}leather\textquotedblright{}, is less likely.
\end{cousinweb}

\begin{grammarlens}[title={What you've built}]
One more word for this chapter.
\end{grammarlens}

\subsection*{Your turn}
\begin{itemize}
\raggedright
  \item \emph{Say it:} \emph{scūtum, scūtī}
  \item \emph{Say it:} \emph{scūtum, scūtī}, once more
  \item \emph{Say it:} say \emph{scūtum, scūtī}
  \item \emph{Recall:} say the Latin for a club, then the Latin for a chain, then say \emph{scūtum, scūtī} again
\end{itemize}

\begin{tcolorbox}[breakable,colback=teal!4,colframe=teal!35!black,title={Before you move on}]
What does scūtum, scūtī mean? (\textquotedblleft{}A shield\textquotedblright{}.) Say it once more.
\end{tcolorbox}
